\documentclass[../../main.tex]{subfiles}

\begin{document}

\section{可测函数}\label{section:ch18-measurable-functions}

对一个可测空间 $(X,\mathcal{M})$, $X$ 上的可测函数的概念等同于一元实变量函数关于 Lebesgue 测度可测的概念. 以下命题的证明与实直线上的 Lebesgue 测度相应结果的证明完全相同, 见 3.1 节命题 1.

\begin{proposition}\label{proposition:ch18-equivalent-conditions}
令 $(X,\mathcal{M})$ 为可测空间而 $f$ 是定义在 $X$ 上的扩充实值函数. 则以下陈述等价:
\begin{description}
\item[（i）] 对每个实数 $c$, 集合 $\{x\in X\mid f(x)<c\}$ 是可测的.
\item[（ii）] 对每个实数 $c$, 集合 $\{x\in X\mid f(x)\leqslant c\}$ 是可测的.
\item[（iii）] 对每个实数 $c$, 集合 $\{x\in X\mid f(x)>c\}$ 是可测的.
\item[（iv）] 对每个实数 $c$, 集合 $\{x\in X\mid f(x)\geqslant c\}$ 是可测的.
\end{description}
这些性质的每个蕴涵对每个实数 $c$, 集合 $\{x\in X\mid f(x)=c\}$ 是可测的.
\end{proposition}

\begin{definition}[可测函数]\label{definition:ch18-measurable-function}
令 $(X,\mathcal{M})$ 为可测空间. $X$ 上的扩充实值函数 $f$ 称为是可测的(关于 $\mathcal{M}$ 可测), 若命题\cref{proposition:ch18-equivalent-conditions}的一个陈述成立, 则所有的四个陈述成立.
\end{definition}

对于集合 $X$ 和 $X$ 的所有子集的 $\sigma$ 代数 $\mathcal{M}=2^{X}$, $X$ 上的每个扩充实值函数关于 $\mathcal{M}$ 都是可测的. 在相反的极端情况下, 考虑 $\sigma$ 代数 $\mathcal{M}=\{X,\varnothing\}$, 仅有的可测函数是常函数. 若 $X$ 是拓扑空间, 而 $\mathcal{M}$ 是 $X$ 的包含 $X$ 上的拓扑的 $\sigma$ 代数, 则 $X$ 上的每个连续实值函数关于 $\mathcal{M}$ 都是可测的. 在第一部分, 我们研究了关于 Lebesgue 可测集的 $\sigma$ 代数可测的一元实变量函数.

由于实数的有界开区间是两个无界开区间的交, 且实数的开集是开区间族的可数并, 我们有以下实值可测函数的刻画(也见习题 1).

\begin{proposition}\label{proposition:ch18-open-preimage}
令 $(X,\mathcal{M})$ 为可测空间, 而 $f$ 是定义在 $\mathcal{M}$ 上的扩充实值函数. 则 $f$ 是可测的当且仅当对每个实数的开集 $\mathcal{O}$, $f^{-1}(\mathcal{O})$ 是可测的.
\end{proposition}

对一个可测空间 $(X,\mathcal{M})$ 和 $X$ 的可测子集 $E$, 我们称定义在 $E$ 上的扩充实值函数 $f$ 为可测的, 若它关于可测空间 $(E,\mathcal{M}_{E})$ 是可测的, 其中 $\mathcal{M}_{E}$ 是 $\mathcal{M}$ 中包含于 $E$ 的集合所组成的族. $X$ 上的可测函数在可测集上的限制是可测的. 此外, 对 $X$ 的扩充实值函数 $f$ 与 $X$ 的可测子集 $E$, $f$ 在 $E$ 和 $X\sim E$ 上的限制都是可测的当且仅当 $f$ 在 $X$ 上是可测的.

\begin{proposition}\label{proposition:ch18-restriction}
令 $(X,\mathcal{M},\mu)$ 为完备的测度空间, 而 $X_{0}$ 是 $X$ 的可测子集. 则 $X$ 的扩充实值函数 $f$ 是可测的当且仅当它在 $X_{0}$ 上的限制是可测的. 特别地, 若 $g$ 和 $h$ 是 $X$ 上的扩充实值函数, 满足在 $X$ 上几乎处处 $g=h$, 则 $g$ 是可测的当且仅当 $h$ 是可测的.
\end{proposition}

\begin{proof}
定义 $f_{0}$ 为 $f$ 在 $X_{0}$ 上的限制. 令 $c$ 为实数而 $E=(c,\infty)$. 若 $f$ 是可测的, 则 $f^{-1}(E)$ 是可测的, 因此 $f^{-1}(E)\cap X_{0}=f_{0}^{-1}(E)$ 也如此. 所以 $f_{0}$ 是可测的. 现在假设 $f_{0}$ 是可测的. 则
\begin{align*}
f^{-1}(E)=f_{0}^{-1}(E)\cup A,
\end{align*}
其中 $A$ 是 $X\sim X_{0}$ 的子集. 由于 $(X,\mathcal{M},\mu)$ 是完备的, $A$ 是可测的, 因此 $f^{-1}(E)$ 也如此. 因此函数 $f$ 是可测的. 第二个断言从第一个得到.
\end{proof}

若测度空间 $(X,\mathcal{M},\mu)$ 不是完备的, 则该命题不成立(见习题 2). 以下定理的证明与实直线上的 Lebesgue 测度的情形的证明相同, 见 3.1 节定理 6.

\begin{theorem}\label{theorem:ch18-algebra-measurable}
令 $(X,\mathcal{M})$ 为可测空间, 而 $f$ 和 $g$ 是 $X$ 上的可测实值函数.
\begin{description}
\item[（线性）] 对任何实数 $\alpha$ 和 $\beta$, $\alpha f+\beta g$ 是可测的.
\item[（乘积）] $f\cdot g$ 是可测的.
\item[（最大与最小值）] 函数 $\max\{f,g\}$ 与 $\min\{f,g\}$ 是可测的.
\end{description}
\end{theorem}

\begin{remark}
两个扩充的实值函数的和在使得函数取符号相反的无穷值的那些点没有定义. 尽管如此, 在可积函数的线性空间的研究中有必要考虑扩充实值可测函数的线性组合. 对几乎处处有限的可测函数, 如同我们对一元实变量函数所做的那样处理. 事实上, 对一个测度空间 $(X,\mathcal{M},\mu)$, 考虑 $X$ 上的两个几乎处处有限的扩充实值可测函数 $f$ 和 $g$. 定义 $X_{0}$ 为 $X$ 中使得 $f$ 和 $g$ 都有限的点集. 由于 $f$ 和 $g$ 是可测函数, $X_{0}$ 是可测集. 此外, $\mu(X\sim X_{0})=0$. 对实数 $\alpha$ 和 $\beta$, 线性组合 $\alpha f+\beta g$ 是 $X_{0}$ 上恰当定义的实值函数. 我们说 $\alpha f+\beta g$ 在 $X$ 上是可测的, 若它在 $X_{0}$ 上的限制关于可测空间 $(X_{0},\mathcal{M}_{0})$ 是可测的, 其中 $\mathcal{M}_{0}$ 是 $\mathcal{M}$ 中所有包含于 $X_{0}$ 的子集的 $\sigma$ 代数. 若 $(X,\mathcal{M},\mu)$ 是完备的, 命题\cref{proposition:ch18-restriction}告诉我们, 该定义等价于在 $\alpha f+\beta g$ 中的一个, 因此任何到整个 $X$ 上的延拓是 $X$ 上的扩充实值可测函数. 我们将 $X$ 上的函数 $\alpha f+\beta g$ 视为 $X$ 上任何与 $\alpha f+\beta g$ 在 $X_{0}$ 上一致的可测的扩充实值函数. 类似的考虑适用于 $f$ 和 $g$ 的乘积与它们的最大值与最小值函数. 有了这个约定, 若扩充的实值可测函数 $f$ 和 $g$ 在 $X$ 上 a.e. 有限, 则前一个定理成立.
\end{remark}

我们已经看到实变量 Lebesgue 可测函数的复合不是可测的(见 3.1 节的例子). 然而, 以下的复合准则非常有用. 它告诉我们若 $f$ 是可测函数而 $0<p<\infty$, 则 $|f|^{p}$ 也是可测的.

\begin{proposition}\label{proposition:ch18-composition}
令 $(X,\mathcal{M})$ 为可测空间, $f$ 是 $X$ 上的可测实值函数, 而 $\varphi:\mathbf{R}\to\mathbf{R}$ 是连续的. 则复合 $\varphi\circ f:X\to\mathbf{R}$ 也是可测的.
\end{proposition}

\begin{proof}
令 $\mathcal{O}$ 为实数的开集. 由于 $\varphi$ 是连续的, $\varphi^{-1}(\mathcal{O})$ 是开的. 因此, 根据命题\cref{proposition:ch18-open-preimage}, $f^{-1}(\varphi^{-1}(\mathcal{O}))=(\varphi\circ f)^{-1}(\mathcal{O})$ 是可测集, 所以 $\varphi\circ f$ 是可测函数.
\end{proof}

如同一元实变量 Lebesgue 可测函数的特殊情形, 可测函数的一个基本的重要性质——函数的可测性——在逐点的极限是保持的.

\begin{theorem}\label{theorem:ch18-limit-measurable}
令 $(X,\mathcal{M},\mu)$ 为测度空间, 而 $\{f_{n}\}$ 是 $X$ 上的可测函数序列, 使得在 $X$ 上 a.e. 逐点地 $\{f_{n}\}\to f$. 若测度空间 $(X,\mathcal{M},\mu)$ 是完备的或收敛在整个 $X$ 上是逐点的, 则 $f$ 是可测的.
\end{theorem}

\begin{proof}
根据命题\cref{proposition:ch18-restriction}, 通过可能地从 $X$ 中去掉测度为零的集合, 我们假定序列 $\{f_{n}\}$ 在整个 $X$ 上逐点收敛. 固定实数 $c$. 我们必须证明集合 $\{x\in X\mid f(x)<c\}$ 是可测的. 观察到对于点 $x\in X$, 由于 $\lim_{n\to\infty}f_{n}(x)=f(x)$, $f(x)<c$ 当且仅当存在自然数 $n$ 和 $k$ 使得对所有 $j\geqslant k$, $f_{j}(x)<c-1/n$. 但对任何自然数 $n$ 和 $j$, 由于函数 $f_{j}$ 是可测的, 集合 $\{x\in X\mid f_{j}(x)<c-1/n\}$ 是可测的. 由于 $\mathcal{M}$ 关于可数交封闭, 对任何 $k$,
\begin{align*}
\bigcap_{j=k}^{\infty}\{x\in X\mid f_{j}(x)<c-1/n\}
\end{align*}
也是可测的. 因此,
\begin{align*}
\{x\in X\mid f(x)<c\}=\bigcup_{1\leqslant k,n<\infty}\left[\bigcap_{j=k}^{\infty}\{x\in X\mid f_{j}(x)<c-1/n\}\right]
\end{align*}
是可测的, 由于 $\mathcal{M}$ 关于可数并封闭.
\end{proof}

若测度空间不是完备的, 该定理不成立(见习题 3).

\begin{corollary}\label{corollary:ch18-limsup-measurable}
令 $(X,\mathcal{M},\mu)$ 为测度空间, 而 $\{f_{n}\}$ 是 $X$ 上的可测函数序列. 则以下函数
\begin{align*}
\sup\{f_{n}\},\quad\inf\{f_{n}\},\quad\limsup\{f_{n}\},\quad\liminf\{f_{n}\}
\end{align*}
是可测的.
\end{corollary}

\begin{definition}[简单函数]\label{definition:ch18-simple-function}
令 $(X,\mathcal{M})$ 为可测空间. 对可测集 $E$, 它的特征函数 $\chi_{E}$ 是在 $E$ 上取值 $1$ 而在 $X\sim E$ 上取值 $0$ 的函数. $X$ 上的实值函数 $\psi$ 称为是简单的, 若存在可测集族 $\{E_{k}\}_{k=1}^{n}$ 和相应的实数集 $\{c_{k}\}_{k=1}^{n}$ 使得
\begin{align*}
\text{在 }X\text{ 上},\quad\psi=\sum_{k=1}^{n}c_{k}\cdot\chi_{E_{k}}.
\end{align*}
观察到 $X$ 上的简单函数是 $X$ 上取有限个实数值的可测实值函数.
\end{definition}

\begin{lemma}[简单逼近引理]\label{lemma:ch18-simple-approximation}
令 $(X,\mathcal{M})$ 为可测空间而 $f$ 是 $X$ 上的有界可测函数, 即存在 $M\geqslant 0$ 使得在 $X$ 上 $|f|\leqslant M$. 则对每个 $\varepsilon>0$, 存在定义在 $X$ 上的简单函数 $\varphi_{\varepsilon}$ 与 $\psi_{\varepsilon}$, 它们具有以下逼近性质:
\begin{align*}
\text{在 }X\text{ 上},\quad\varphi_{\varepsilon}\leqslant f\leqslant\psi_{\varepsilon}\text{ 且 }0\leqslant\psi_{\varepsilon}-\varphi_{\varepsilon}<\varepsilon.
\end{align*}
\end{lemma}

\begin{proof}
令 $[c,d]$ 为包含 $X$ 的象 $f(X)$ 的有界区间, 而
\begin{align*}
c=y_{0}<y_{1}<\cdots<y_{n-1}<y_{n}=d
\end{align*}
是闭有界区间 $[c,d]$ 的一个分划, 使得对 $1\leqslant k\leqslant n$, $y_{k}-y_{k-1}<\varepsilon$. 对 $1\leqslant k\leqslant n$ 定义
\begin{align*}
I_{k}=(y_{k-1},y_{k})\text{ 与 }X_{k}=f^{-1}(I_{k}).
\end{align*}
由于每个 $I_{k}$ 是开区间而函数 $f$ 是可测的, 每个集合 $X_{k}$ 是可测的. 定义 $X$ 上的简单函数 $\varphi_{\varepsilon}$ 和 $\psi_{\varepsilon}$ 为
\begin{align*}
\varphi_{\varepsilon}=\sum_{k=1}^{n}y_{k-1}\cdot\chi_{X_{k}}\text{ 与 }\psi_{\varepsilon}=\sum_{k=1}^{n}y_{k}\cdot\chi_{X_{k}}.
\end{align*}
令 $x$ 属于 $X$. 由于 $f(X)\subseteq[c,d]$, 存在唯一的 $k$, 使得 $y_{k-1}\leqslant f(x)<y_{k}$. 因此,
\begin{align*}
\varphi_{\varepsilon}(x)=y_{k-1}\leqslant f(x)<y_{k}=\psi_{\varepsilon}(x).
\end{align*}
但 $y_{k}-y_{k-1}<\varepsilon$, 因此 $\varphi_{\varepsilon}$ 与 $\psi_{\varepsilon}$ 有所要的逼近性质.
\end{proof}

\begin{theorem}[简单逼近定理]\label{theorem:ch18-simple-approximation-thm}
令 $(X,\mathcal{M},\mu)$ 为测度空间而 $f$ 是 $X$ 上的可测函数. 存在定义在 $X$ 上的简单函数序列 $\{\psi_{n}\}$, 它在 $X$ 上逐点收敛于 $f$ 且具有性质:
\begin{align*}
\text{在 }X\text{ 上},\quad\text{对所有 }n,\quad|\psi_{n}|\leqslant|f|.
\end{align*}
\begin{description}
\item[（i）] 若 $X$ 是 $\sigma$ 有限的, 则我们可以选取序列 $\{\psi_{n}\}$ 使得每个 $\psi_{n}$ 在有限测度集外消失.
\item[（ii）] 若 $f$ 是非负的, 我们可以选取序列 $\{\psi_{n}\}$ 为递增的且在 $X$ 上每个 $\psi_{n}\geqslant 0$.
\end{description}
\end{theorem}

\begin{proof}
固定一个自然数 $n$. 定义 $E_{n}=\{x\in X\mid|f(x)|\leqslant n\}$. 由于 $|f|$ 是可测函数, $E_{n}$ 是可测子集, 而 $f$ 在 $E_{n}$ 上的限制是有界可测函数. 将简单逼近引理用于 $f$ 在 $E_{n}$ 上的限制, 其中选取 $\varepsilon=1/n$, 我们可以选取 $E_{n}$ 上的简单函数 $h_{n}$ 与 $g_{n}$, 它们具有以下逼近性质:
\begin{align*}
\text{在 }E_{n}\text{ 上},\quad h_{n}\leqslant f\leqslant g_{n}\text{ 且 }0\leqslant g_{n}-h_{n}<1/n.
\end{align*}
对 $E_{n}$ 中的 $x$, 若 $f(x)=0$, 定义 $\psi_{n}(x)=0$; 若 $f(x)>0$, $\psi_{n}(x)=\max\{h_{n}(x),0\}$; 若 $f(x)<0$, $\psi_{n}(x)=\min\{g_{n}(x),0\}$. 通过定义若 $f(x)>n$ 令 $\psi_{n}(x)=n$, 若 $f(x)<-n$ 令 $\psi_{n}(x)=-n$, 将 $\psi_{n}$ 延拓到整个 $X$ 上. 这定义了 $X$ 上的简单函数序列 $\{\psi_{n}\}$. 如同对实变量 Lebesgue 可测函数相应结果证明所做的那样(见 3.2 节), 对每个 $n$, 在 $X$ 上, $|\psi_{n}|\leqslant|f|$ 且序列 $\{\psi_{n}\}$ 在 $X$ 上逐点收敛于 $f$.

若 $X$ 是 $\sigma$ 有限的, 将 $X$ 表示为可测子集的可数上升族 $\{X_{n}\}_{n=1}^{\infty}$ 的并, 每个具有有限测度. 用 $\psi_{n}\cdot\chi_{X_{n}}$ 代替每个 $\psi_{n}$, 从而(i)得证. 若 $f$ 非负, 用 $\max_{1\leqslant i\leqslant n}|\psi_{i}|$ 代替每个 $\psi_{n}$, 从而(ii)得证.
\end{proof}

如同一元实变量 Lebesgue 可测函数的情形, 以下一般形式的 Egoroff 定理的证明从测度的连续性与可数可加性得到, 见 3.3 节.

\begin{theorem}[Egoroff]\label{theorem:ch18-egoroff}
令 $(X,\mathcal{M},\mu)$ 为有限测度空间, 而 $\{f_{n}\}$ 是 $X$ 上的可测函数序列, 它在 $X$ 上 a.e. 逐点收敛于在 $X$ 上 a.e. 有限的函数 $f$. 则对每个 $\varepsilon>0$, 存在 $X$ 的可测子集 $X_{\varepsilon}$ 使得
\begin{align*}
\text{在 }X_{\varepsilon}\text{ 上一致地 }\{f_{n}\}\to f,\quad\mu(X\sim X_{\varepsilon})<\varepsilon.
\end{align*}
\end{theorem}

\begin{exercise}\label{exercise:ch18-01}
证明 $X$ 上的扩充实值函数 $f$ 是可测的当且仅当 $f^{-1}\{\infty\}$ 和 $f^{-1}\{-\infty\}$ 是可测的, 从而对实数的每个 Borel 集 $E$, $f^{-1}(E)$ 也可测.
\end{exercise}

\begin{exercise}\label{exercise:ch18-02}
假定 $(X,\mathcal{M},\mu)$ 不是完备的. 令 $E$ 为不属于 $\mathcal{M}$ 的测度为零的子集. 在 $X$ 上令 $f=0$ 而 $g=\chi_{E}$. 证明在 $X$ 上 $f=g$ a.e., 但 $f$ 是可测的而 $g$ 不是.
\end{exercise}

\begin{exercise}\label{exercise:ch18-03}
假定 $(X,\mathcal{M},\mu)$ 不是完备的. 证明存在 $X$ 上的可测函数序列 $\{f_{n}\}$, 它在 $X$ 上 a.e. 逐点收敛于不可测函数 $f$.
\end{exercise}

\begin{exercise}\label{exercise:ch18-04}
令 $E$ 为 $X$ 的可测子集而 $f$ 是 $X$ 上的扩充实值函数. 证明 $f$ 是可测的当且仅当它在 $E$ 和 $X\sim E$ 上的限制都是可测的.
\end{exercise}

\begin{exercise}\label{exercise:ch18-05}
证明 $X$ 上的扩充实值函数 $f$ 是可测的当且仅当对每个有理数 $c$, $\{x\in X\mid f(x)<c\}$ 是可测集.
\end{exercise}

\begin{exercise}\label{exercise:ch18-06}
考虑 $X$ 上的两个在 $X$ 上 a.e. 有限的扩充实值函数 $f$ 和 $g$. 定义 $X_{0}$ 为 $X$ 的使得 $f$ 和 $g$ 都有限的点集. 证明 $X_{0}$ 是可测的且 $\mu(X\sim X_{0})=0$.
\end{exercise}

\begin{exercise}\label{exercise:ch18-07}
令 $X$ 为非空集. 证明 $X$ 上的每个扩充实值函数关于可测空间 $(X,2^{X})$ 是可测的.
\begin{description}
\item[（i）] 令 $x_{0}$ 属于 $X$ 而 $\delta_{x_{0}}$ 是 $2^{X}$ 上在 $x_{0}$ 的 Dirac 测度. 证明 $X$ 的两个函数 a.e. $[\delta_{x_{0}}]$ 相等当且仅当它们在 $x_{0}$ 取同样的值.
\item[（ii）] 令 $\eta$ 为 $2^{X}$ 上的计数测度. 证明 $X$ 上的两个函数 a.e. $[\eta]$ 相等当且仅当它们在 $X$ 中的每个点取同样的值.
\end{description}
\end{exercise}

\begin{exercise}\label{exercise:ch18-08}
令 $X$ 为拓扑空间, 而 $\mathcal{B}(X)$ 是包含 $X$ 上的拓扑的最小 $\sigma$ 代数. $\mathcal{B}(X)$ 称为与拓扑空间 $X$ 伴随的 Borel $\sigma$ 代数. 证明 $X$ 上的任何连续实值函数关于 Borel 可测空间 $(X,\mathcal{B}(X))$ 是可测的.
\end{exercise}

\begin{exercise}\label{exercise:ch18-09}
若 $\mathbf{R}$ 上的实值函数关于 Lebesgue 可测集的 $\sigma$ 代数可测, 它必须关于 Borel 可测空间 $(\mathbf{R},\mathcal{B}(\mathbf{R}))$ 可测吗?
\end{exercise}

\begin{exercise}\label{exercise:ch18-10}
验证命题\cref{proposition:ch18-equivalent-conditions}和定理\cref{theorem:ch18-algebra-measurable}的证明可从实直线上的 Lebesgue 测度的相应结果的证明得出.
\end{exercise}

\begin{exercise}\label{exercise:ch18-11}
证明推论\cref{corollary:ch18-limsup-measurable}.
\end{exercise}

\begin{exercise}\label{exercise:ch18-12}
证明 Egoroff 定理. 若不假设极限函数是 a.e. 有限的, Egoroff 定理成立吗?
\end{exercise}

\begin{exercise}\label{exercise:ch18-13}
令 $\{f_{n}\}$ 为 $X$ 上的实值可测函数序列, 使得对每个自然数 $n$, $\mu\{x\in X\mid|f_{n}(x)-f_{n+1}(x)|>1/2^{n}\}<1/2^{n}$. 证明在 $X$ 上 $\{f_{n}\}$ a.e. 逐点收敛于 $f$.（提示: 用 Borel-Cantelli 引理.）
\end{exercise}

\begin{exercise}\label{exercise:ch18-14}
在 Egoroff 定理的假设下, 证明 $X=\bigcup_{k=0}^{\infty}X_{k}$, 其中每个 $X_{k}$ 是可测的, $\mu(X_{0})=0$ 且对 $k\geqslant 1$ 在 $X_{k}$ 上 $\{f_{n}\}$ 一致收敛到 $f$.
\end{exercise}

\begin{exercise}\label{exercise:ch18-15}
$X$ 上的可测实值函数序列 $\{f_{n}\}$ 称为依测度收敛到可测函数 $f$, 若对每个 $\eta>0$,
\begin{align*}
\lim_{n\to\infty}\mu\{x\in X\mid|f_{n}(x)-f(x)|>\eta\}=0.
\end{align*}
可测函数序列 $\{f_{n}\}$ 称为是依测度 Cauchy 的, 若对每个 $\varepsilon>0$ 与 $\eta>0$, 存在指标 $N$ 使得对每个 $m,n\geqslant N$,
\begin{align*}
\mu\{x\in X\mid|f_{n}(x)-f_{m}(x)|>\eta\}<\varepsilon.
\end{align*}
\begin{description}
\item[（i）] 证明: 若 $\mu(X)<\infty$ 且 $\{f_{n}\}$ 在 $X$ 上 a.e. 逐点收敛到可测函数 $f$, 则 $\{f_{n}\}$ 依测度收敛到 $f$.（提示: 用 Egoroff 定理.）
\item[（ii）] 证明: 若 $\{f_{n}\}$ 依测度收敛到 $f$, 则存在 $\{f_{n}\}$ 的子序列在 $X$ 上 a.e. 逐点收敛到 $f$.（提示: 用 Borel-Cantelli 引理.）
\item[（iii）] 证明: 若 $\{f_{n}\}$ 是依测度 Cauchy 的, 则存在可测函数 $f$ 使得 $\{f_{n}\}$ 依测度收敛到 $f$.
\end{description}
\end{exercise}

\begin{exercise}\label{exercise:ch18-16}
假设 $\mu(X)<\infty$. 证明依测度 $\{f_{n}\}\to f$ 当且仅当 $\{f_{n}\}$ 的每个子序列有一个子序列在 $X$ 上 a.e. 逐点收敛于 $f$. 用此证明对两个依测度收敛的序列, 乘积序列也依测度收敛到极限的乘积.
\end{exercise}

\end{document}